\documentclass[12pt, a4paper]{article}

% ── Packages ─────────────────────────────────────────────────────────────────
\usepackage[utf8]{inputenc}
\usepackage[T1]{fontenc}
\usepackage{lmodern}
\usepackage[margin=1in]{geometry}
\usepackage{graphicx}
\usepackage{booktabs}
\usepackage{hyperref}
\usepackage{amsmath}
\usepackage{enumitem}
\usepackage{titlesec}
\usepackage{fancyhdr}
\usepackage{xcolor}
\usepackage{float}
\usepackage{caption}
\usepackage{listings}

\hypersetup{
    colorlinks=true,
    linkcolor=blue!70!black,
    citecolor=blue!70!black,
    urlcolor=blue!70!black
}

% ── Header / Footer ─────────────────────────────────────────────────────────
\pagestyle{fancy}
\fancyhf{}
\rhead{Final Report — M1 \& M2}
\lhead{Intelligent News Credibility Analyzer}
\cfoot{\thepage}

% ── Title ────────────────────────────────────────────────────────────────────
\title{
    \vspace{-1cm}
    \textbf{\LARGE Intelligent News Credibility Analyzer} \\[0.4cm]
    \large From Classical Machine Learning to Agentic GenAI Fact-Checking \\[0.3cm]
    \large End-Semester Capstone Project Report
}

\author{
    Aditya Mathur \and Om Kar Shukla \and Yachna Khanna
}

\date{April 2026}

\begin{document}

\maketitle
\thispagestyle{fancy}

% ══════════════════════════════════════════════════════════════════════════════
\section{Abstract}
% ══════════════════════════════════════════════════════════════════════════════

This capstone project tackles the automated detection of fake news utilizing two contrasting artificial intelligence paradigms. \textbf{Milestone 1} focused on establishing a robust statistical baseline using classical Machine Learning. By extracting TF-IDF features from a dataset of 38,644 articles and applying a Logistic Regression classifier, the system achieved a 96.7\% accuracy. Though performant on its training distribution, the ML model lacks the capability to verify novel factual assertions. 

\textbf{Milestone 2} addresses this fundamental limitation by introducing an \textbf{Agentic AI framework} powered by Large Language Models (LLMs). Utilizing Google's Llama 3 70B within a LangChain ReAct (Reasoning and Acting) orchestration loop, the agent autonomously retrieves real-world information using DuckDuckGo web search. The architecture leverages Retrieval-Augmented Generation (RAG) principles to synthesize live web evidence into a conclusive, interpretable credibility verdict complete with confidence metrics and cited sources. Both phases are successfully integrated into an interactive, publicly deployed Streamlit application, yielding a comprehensive and modern pipeline for misinformation detection.

\vspace{0.3cm}
\noindent\textbf{Repository:} \url{https://github.com/adityamathur5836/News-Credibility-Analyzer}

% ══════════════════════════════════════════════════════════════════════════════
\section{Milestone 1: Classical Machine Learning Baseline}
% ══════════════════════════════════════════════════════════════════════════════

Milestone 1 established our statistical baseline using supervised machine learning. 

\subsection{Data Processing \& Inference Pipeline}
The system was trained on a balanced corpus derived from two standard datasets (\texttt{Fake.csv} and \texttt{True.csv}), resulting in 38,644 unique, labeled articles following rigorous deduplication and null-handling. The text data was subjected to a classical NLP pipeline (lowercasing, punctuation stripping, tokenization via NLTK, stopword omission, and lemmatization) before being vectorized utilizing a Term Frequency-Inverse Document Frequency (TF-IDF) approach constrained to a 5,000-word vocabulary.

A tuned Logistic Regression classifier was identified as the optimal inference mechanism. Its probability distributions facilitate granular confidence scoring within the user interface, improving upon the binary rigidity exhibited by Decision Tree architectures tested in parallel.

\subsection{Milestone 1 Performance}
Evaluated on a randomized test subset, the Logistic Regression model attained an \textbf{Accuracy of 96.70\%} and an \textbf{F1 Score of 96.59\%}. Despite this high empirical performance, the model primarily evaluates linguistic pattern congruence rather than factual accuracy, establishing the necessity for the Agentic pipeline developed in Phase 2.

% ══════════════════════════════════════════════════════════════════════════════
\section{Milestone 2: Agentic Workflow \& Large Language Models}
% ══════════════════════════════════════════════════════════════════════════════

Milestone 2 enhances the application via a multi-agent architectural loop designed not to simply classify text, but to investigate it. 

\subsection{ReAct Agent Framework}
The core intelligence engine leverages the \textbf{ReAct} (Reasoning + Acting) capability within the LangChain framework. The ReAct prompt instructs the LLM (Groq Llama 3 1.5 Flash) to articulate its chain of thought, formulate a strategy, invoke tools, and observe tool outputs until it converges upon a final answer. The ReAct trajectory for a typical query follows:
\begin{enumerate}
    \item \textbf{Question:} "Is it true that [Claim X] occurred?"
    \item \textbf{Thought:} I need to first run the ML pre-screener, then retrieve live evidence if the claim demands verification.
    \item \textbf{Action:} Execute Web Search on [Keywords]
    \item \textbf{Observation:} [DuckDuckGo Snippets]
    \item \textbf{Final Answer:} Synthesis of the retrieved truth with appended source URLs.
\end{enumerate}

\subsection{Tool Integration}
The agent is bestowed with two distinct parametric capabilities:
\begin{enumerate}
    \item \texttt{ml\_prescreener}: Invokes the local Logistic Regression model to acquire statistical probabilities, grounding the agent's initial bias.
    \item \texttt{web\_search}: Implements the \texttt{duckduckgo-search} Python library to scrape real-time internet search snippets, mimicking a human researcher fact-checking an assertion.
\end{enumerate}

% ══════════════════════════════════════════════════════════════════════════════
\section{Retrieval-Augmented verification}
% ══════════════════════════════════════════════════════════════════════════════

While an off-the-shelf LLM is prone to hallucination—particularly with novel or obscure political claims—our architecture mitigates this via Retrieval-Augmented Generation (RAG). By forcing the agent to extract semantic facts from internet queries \textit{before} generating an output, we anchor the LLM's vast parametric knowledge in empirical, up-to-date reality. 

\subsection{Design Justification}
\begin{itemize}
    \item \textbf{Groq Llama 3 1.5 Flash}: Chosen as the underlying LLM for its exceptional inference speed, expansive context window, and accessibility under free-tier API constraints, satisfying the budgetary requirements of the project.
    \item \textbf{LangChain Orchestration}: Allows seamless integration of bespoke tools (like our bespoke ML classifier) into standardized LLM prompt chains.
    \item \textbf{FAISS / Local Embeddings (Modular)}: For multi-document contexts, the architecture trivially scales by utilizing CPU-bound \texttt{sentence-transformers} traversing an in-memory FAISS indices, ensuring efficient inner-product matching.
\end{itemize}

% ══════════════════════════════════════════════════════════════════════════════
\section{Qualitative Analysis \& Hosted Deployment}
% ══════════════════════════════════════════════════════════════════════════════

\subsection{Case Observation}
When fed a fabricated recent headline—such as claims regarding fictitious planetary alignments or nonexistent legislative acts—the Milestone 1 ML model often expresses high confidence (due to matching structural linguistic patterns). Conversely, the Milestone 2 Agent invokes the Web Search tool, observes a total lack of credible journalistic reporting corroborating the headline, and correctly labels the claim as \textit{Fake/Suspicious}, providing the thought process and null searches as logic.

\subsection{Cloud Deployment \& Usability}
The finalized artifact is a production-ready application deployed on \textbf{Streamlit Community Cloud}. The User Interface incorporates a dedicated tab-layout enabling users to juxtapose the rapid, statistical outputs of Milestone 1 against the thorough, epistemological investigative trace produced by the Milestone 2 GenAI Agent. The deployment securely manages the Groq API keys via Streamlit's encrypted \texttt{secrets.toml} backend, ensuring security standards aligned with professional engineering best practices.

\end{document}
